\documentclass[10pt]{scrartcl}
\usepackage[mdthm,colorsec,noauthor,noasy]{evan}
\usepackage{handout}
\DeclareMathOperator{\im}{range}
\title{MATH 110 --- Week 5}
\subtitle{Notes Chapters 9--10; Homework 5}
\author{}\date{}\ihead{}
\begin{document}\maketitle
\section{Quotient maps and the first isomorphism theorem}\begin{qbox}Which information about a vector is determined by its image under a linear map? For $T(x,y)=x$, inputs with the same first coordinate have the same image; their differences lie in $\ker T=\mathbb R(0,1)$.\end{qbox}
\begin{prob}\begin{definition*}[Notes Definition 25; LADR 3.104]For a subspace $U\subseteq V$, the quotient map is $\pi:V\to V/U$, $\pi(v)=v+U$. Week 4 proved that cosets form a vector space, with zero coset $U$, and that $v+U=w+U$ iff $v-w\in U$.\end{definition*}\end{prob}
\begin{proposition*}[Notes Definition 25; LADR 3.104]The quotient map is linear and surjective, with kernel $U$.\end{proposition*}
\begin{motive}\ii[Want.] Relate calculations in a quotient to a familiar map on $V$.\ii[Know.] The quotient operations were defined using representative sums and scalar multiples.\ii[Key step.] Apply those definitions and the test for the zero coset.\end{motive}
\begin{prob}\begin{proof}\begin{steps}\ii $\pi(av+w)=(av+w)+U=a(v+U)+(w+U)=a\pi(v)+\pi(w)$. Every coset is $\pi(v)$ for some $v$, giving surjectivity.\ii $v\in\ker\pi$ iff $v+U=U$, which by Week 4 Proposition 39 is equivalent to $v\in U$.\end{steps}\end{proof}\end{prob}
\begin{proposition*}[Notes Proposition 40; LADR 3.105]If $V$ is finite-dimensional, $\dim(V/U)=\dim V-\dim U$.\end{proposition*}
\begin{motive}\ii[Want.] Count the independent cosets left after quotienting.\ii[Know.] The quotient map has range $V/U$ and kernel $U$.\ii[Key step.] Apply rank--nullity to this surjective map.\end{motive}
\begin{prob}\begin{proof}\begin{steps}\ii Week 3 rank--nullity gives $\dim V=\dim\ker\pi+\dim\im\pi$.\ii Substitute $\ker\pi=U$ and $\im\pi=V/U$, then rearrange. The range is finite-dimensional because it is the image of a finite-dimensional domain.\end{steps}\end{proof}\end{prob}
\method{Construct a quotient basis}\begin{steps}\ii Extend a basis $u_1,\ldots,u_k$ of $U$ to a basis $u_1,\ldots,u_k,v_1,\ldots,v_r$ of $V$.\ii The cosets $v_1+U,\ldots,v_r+U$ span: in a basis expansion, the $u_i$ terms belong to the zero coset.\ii If $\sum a_j(v_j+U)=U$, then $\sum a_jv_j\in U$, producing a relation in the full basis; all $a_j=0$. Thus the cosets form a basis.\ii For $U=\operatorname{span}(e_1,e_2)\subseteq\mathbb R^3$, the single coset $e_3+U$ is a quotient basis.\end{steps}
\begin{proposition*}[Notes Proposition 41; LADR 3.106--3.107]For $T:V\to W$ linear, the rule $\overline T:V/\ker T\to\im T$, $\overline T(v+\ker T)=Tv$, is an isomorphism. If $\iota:\im T\to W$ is the inclusion $w\mapsto w$, then $T=\iota\overline T\pi$.\end{proposition*}
\begin{motive}\ii[Want.] Identify exactly the input distinctions that the map preserves.\ii[Know.] $Tv=Tw$ iff $v-w\in\ker T$, and invertibility is equivalent to linear bijectivity.\ii[Key step.] Use a coset as input, making the equal-output inputs one element of the quotient.\end{motive}
\begin{prob}\begin{proof}\begin{steps}\ii If $v+\ker T=w+\ker T$, then $v-w\in\ker T$, so $Tv-Tw=T(v-w)=0$. Thus the rule is well-defined.\ii Using representatives, $\overline T(a(v+\ker T)+(w+\ker T))=T(av+w)=aTv+Tw$, proving linearity.\ii If $\overline T(v+\ker T)=0$, then $v\in\ker T$, so the input coset is zero. Thus the map is injective. It is onto by the definition of $\im T$.\ii Week 4 Proposition 29 makes it an isomorphism. For each $v$, $\iota\overline T\pi(v)=Tv$, proving the factorization.\end{steps}\end{proof}\end{prob}
\begin{center}\begin{tikzpicture}[>=stealth]\node (v)at(0,1.5){$V$};\node(w)at(5,1.5){$W$};\node(q)at(0,0){$V/\ker T$};\node(r)at(5,0){$\im T$};\draw[->](v)--node[above]{$T$}(w);\draw[->](v)--node[left]{$\pi$}(q);\draw[->](q)--node[below]{$\overline T$}(r);\draw[->](r)--node[right]{$\iota$}(w);\end{tikzpicture}\\Both routes from $V$ to $W$ give the same image $Tv$.\end{center}
\begin{example*}[Notes Example 43]Differentiation $D:\mathcal P(\mathbb R)\to\mathcal P(\mathbb R)$ is onto and has kernel the constants $\mathbb R$, as proved in Week 4.
\begin{steps}\ii Proposition 41 gives $\overline D:\mathcal P(\mathbb R)/\mathbb R\to\mathcal P(\mathbb R)$, $p+\mathbb R\mapsto p^{\prime}$.
\ii Its inverse sends $q$ to the coset $r+\mathbb R$ of any antiderivative $r$ of $q$. Such an $r$ exists by termwise integration; any two differ by a constant, so the coset is unique.
\ii For $q=x^2$, the inverse image is $\frac13x^3+\mathbb R$. Thus the indefinite integral denotes one quotient vector rather than one chosen polynomial.\end{steps}\end{example*}
\begin{prob}\textbf{Maps and subspaces.}\begin{steps}\ii An injective $T:V\to W$ gives an isomorphism $V\to\im T$ by restricting its codomain; it is onto that range by definition.\ii A surjective $T:V\to W$ gives $V/\ker T\cong W$ by Proposition 41.\ii In finite dimension, $\dim(V/\ker T)=\dim V-\dim\ker T=\dim\im T$. This recovers the rank--nullity equality, but is not an independent proof because Proposition 40 used rank--nullity.\end{steps}\end{prob}
\begin{shortcuts}\ii To find a quotient basis, extend a subspace basis and retain the cosets of the added vectors.\ii To identify a quotient with a concrete space, find a surjective map with the specified subspace as kernel; Proposition 41 supplies the isomorphism.\end{shortcuts}\begin{trapbox}\ii The rule $v+U\mapsto Tv$ need not be well-defined unless $U\subseteq\ker T$: with $T=I$ on $\mathbb R^2$ and $U=\mathbb Re_1$, the same coset $U=0+U=e_1+U$ would have different images.\ii Infinite-dimensional $V$ need not give an infinite-dimensional quotient: evaluation $p\mapsto p(0)$ identifies $\mathcal P(\mathbb R)/x\mathcal P(\mathbb R)$ with $\mathbb R$ by Proposition 41.\end{trapbox}
\section{Induced maps and successive quotients}\begin{qbox}When does a linear map send equal cosets to equal cosets? For example, the identity sends $\mathbb Re_1$ into $\mathbb R^2$, so it gives a map from $\mathbb R^2/\mathbb Re_1$ to $\mathbb R^2/\mathbb R^2$.\end{qbox}
\begin{proposition*}[Notes Proposition 42]If $T:V\to W$ is linear, $U_1\subseteq V$ and $U_2\subseteq W$ are subspaces, and $T(U_1)\subseteq U_2$, then $\overline T:V/U_1\to W/U_2$, $\overline T(v+U_1)=Tv+U_2$, is well-defined and linear. This condition is also necessary for that rule to be well-defined.\end{proposition*}
\begin{motive}\ii[Want.] Define a map on quotient vectors without choosing representatives.\ii[Know.] Two representatives differ by a vector in the subspace.\ii[Key step.] Apply $T$ to the difference and use the containment.\end{motive}
\begin{prob}\begin{proof}\begin{steps}\ii If $v+U_1=w+U_1$, then $v-w\in U_1$, hence $Tv-Tw\in U_2$. Thus $Tv+U_2=Tw+U_2$.\ii On representative combinations, $\overline T((av+w)+U_1)=T(av+w)+U_2=a(Tv+U_2)+(Tw+U_2)$, proving linearity.\ii Conversely, if the rule is well-defined, for $u\in U_1$ the equality $u+U_1=0+U_1$ forces $Tu+U_2=U_2$, hence $Tu\in U_2$.\end{steps}\end{proof}\end{prob}
\begin{proposition*}[Notes Proposition 43: second isomorphism theorem]For subspaces $U,W\subseteq V$, $\phi:U/(U\cap W)\to(U+W)/W$, $\phi(u+(U\cap W))=u+W$, is an isomorphism. In finite dimension it gives $\dim U-\dim(U\cap W)=\dim(U+W)-\dim W$.\end{proposition*}
\begin{motive}\ii[Want.] Represent every coset of $W$ in the sum by a vector of $U$.\ii[Know.] The first isomorphism theorem identifies the quotient by the kernel of a surjective map.\ii[Key step.] Use the map $u\mapsto u+W$ and compute its kernel.\end{motive}
\begin{prob}\begin{proof}\begin{steps}\ii Define $R:U\to(U+W)/W$ by $Ru=u+W$. This is linear by quotient operations.\ii Every coset in the target is $(u+w)+W=u+W$, so $R$ is onto.\ii $Ru=W$ iff $u\in W$; since the domain is $U$, $\ker R=U\cap W$. Proposition 41 supplies the stated isomorphism.\ii If the spaces are finite-dimensional, apply Proposition 40 to both quotients and use equal dimensions of isomorphic spaces.\end{steps}\end{proof}\end{prob}
\begin{prob}\textbf{Meaning of the chosen isomorphism.} Here the formula for $\phi$ is specified directly by the subspaces, without choosing bases or complements. This is the sense in which these isomorphisms are called natural in the notes.\end{prob}
\begin{proposition*}[Notes Proposition 44: third isomorphism theorem]For subspaces $W\subseteq U\subseteq V$, $U/W$ is the subspace $\{u+W:u\in U\}$ of $V/W$, and $\phi:(V/W)/(U/W)\to V/U$, $\phi((v+W)+(U/W))=v+U$, is an isomorphism.\end{proposition*}
\begin{motive}\ii[Want.] Compare quotienting in two stages with quotienting directly by $U$.\ii[Know.] Proposition 42 makes representative formulas legitimate, and Proposition 41 identifies their quotient by the kernel.\ii[Key step.] Use the identity map on $V$ to pass from $W$-cosets to $U$-cosets.\end{motive}
\begin{prob}\begin{proof}\begin{steps}\ii The set $U/W$ contains the zero coset and is closed under linear combinations because $U$ is a subspace, so it is a subspace of $V/W$.\ii Since $I(W)\subseteq U$, Proposition 42 gives the linear map $R:V/W\to V/U$, $R(v+W)=v+U$. It is onto because every $v+U$ has this form.\ii $R(v+W)=U$ iff $v\in U$, so $\ker R=U/W$. Proposition 41 gives the displayed formula and isomorphism.\ii For finite-dimensional $V$, its source dimension is $(\dim V-\dim W)-(\dim U-\dim W)=\dim V-\dim U$, agreeing with the target.\end{steps}\end{proof}\end{prob}
\method{Identify successive quotients in coordinates}\begin{steps}\ii In $V=\mathbb R^3$, take $W=\mathbb Re_1$ and $U=\operatorname{span}(e_1,e_2)$. A coset of $W$ is determined by coordinates $(y,z)$: representative differences in $W$ change only $x$.\ii The subspace $U/W$ consists of pairs $(y,0)$. A coset modulo it is determined by $z$.\ii Directly, a coset in $V/U$ is also determined by $z$, and Proposition 44 identifies the two descriptions by keeping that coordinate.\end{steps}
\begin{shortcuts}\ii For an induced map, check $T(U_1)\subseteq U_2$ before doing any calculations.\ii To prove a quotient isomorphism, start from a linear map on the original space and identify its kernel and range.\end{shortcuts}\begin{trapbox}\ii Nested quotients require $W\subseteq U$: for $U=\mathbb Re_1$ and $W=\mathbb Re_2$ in $\mathbb R^2$, the set $W$ is not a subspace of $U$, so the quotient $U/W$ is not defined.\ii Equal dimensions do not prove a specified map is invertible: the zero map $\mathbb R\to\mathbb R$ has equal domain and codomain dimensions but is not an isomorphism.\end{trapbox}
\section{Linear functionals and dual bases}\begin{qbox}How can a vector be described by linear scalar values? On $F^n$, the $j$th coordinate is a linear functional, and all coordinates together determine the vector.\end{qbox}
\begin{prob}\begin{definition*}[Notes Definition 26; LADR 3.108, 3.110]A linear functional on $V$ is a linear map $\varphi:V\to F$. The dual space is $V^{\prime}=L(V,F)$; its elements are also called linear forms, covectors or one-forms. It is a vector space by Week 3 Proposition 18. The prime on $V^{\prime}$ denotes duality, not differentiation of a polynomial.\end{definition*}\end{prob}
\begin{example*}[LADR 3.109: functionals]\begin{steps}\ii On $F^n$, $\varphi(x_1,\ldots,x_n)=\sum_j a_jx_j$ is linear by distributivity, for any fixed scalars $a_j$.
\ii On $\mathcal P(\mathbb R)$, evaluation $p\mapsto p(c)$ and integration $p\mapsto\int_0^1p(t)\,dt$ are linear by coefficient evaluation and the elementary integral linearity cited in Week 3.
\ii If a functional is nonzero, choose $v$ with $\varphi(v)=a\ne0$. For any $b\in F$, $\varphi((b/a)v)=b$, so the functional is onto.\end{steps}\end{example*}
\begin{proposition*}[Notes Chapter 10; LADR 3.111]If $V$ is finite-dimensional, then $\dim V^{\prime}=\dim V$.\end{proposition*}
\begin{motive}\ii[Want.] Count the independent linear scalar functions.\ii[Know.] Week 3 proved $\dim L(V,W)=\dim V\dim W$, and $F$ has basis $(1)$.\ii[Key step.] Use $W=F$ in that formula.\end{motive}
\begin{prob}\begin{proof}\begin{steps}\ii $V^{\prime}=L(V,F)$ and $\dim F=1$.\ii Therefore $\dim V^{\prime}=(\dim V)\cdot1=\dim V$.\end{steps}\end{proof}\end{prob}
\begin{prob}\begin{definition*}[Notes Definition 27; LADR 3.112]For an ordered basis $v_1,\ldots,v_n$ of $V$, its dual list is $\varphi_1,\ldots,\varphi_n\in V^{\prime}$ with $\varphi_j(v_k)=\delta_{jk}$, where $\delta_{jk}=1$ for $j=k$ and $0$ otherwise. Existence and uniqueness of each functional follow from the prescribed-basis-image theorem (Week 3 Proposition 21). We next show this list is a basis, called the dual basis.\end{definition*}\end{prob}
\begin{proposition*}[Notes Proposition 45; LADR 3.114, 3.116]The dual list is a basis of $V^{\prime}$. For every $v\in V$ and $\psi\in V^{\prime}$, $v=\sum_j\varphi_j(v)v_j$ and $\psi=\sum_j\psi(v_j)\varphi_j$.\end{proposition*}
\begin{motive}\ii[Want.] Recover vectors and functionals by evaluating on the appropriate basis.\ii[Know.] The defining values are $\varphi_j(v_k)=\delta_{jk}$, and maps are determined on a basis.\ii[Key step.] Apply one coordinate functional or evaluate a proposed functional identity on one basis vector.\end{motive}
\begin{prob}\begin{proof}\begin{steps}\ii Write $v=\sum a_jv_j$. Applying $\varphi_k$ gives $\varphi_k(v)=a_k$, proving the first expansion.\ii If $\sum b_j\varphi_j=0$, evaluation at $v_k$ gives $b_k=0$. Thus the dual list is independent.\ii The functional $\sum_j\psi(v_j)\varphi_j$ agrees with $\psi$ at every $v_k$, so the two linear maps agree everywhere. This gives the second expansion and proves spanning.\end{steps}\end{proof}\end{prob}
\begin{example*}[Notes Example 44; LADR 3.113]\begin{steps}\ii In the standard basis of $F^n$, a functional has row representation $(a_1\ \cdots\ a_n)$ acting on a coordinate column by $\sum a_jx_j$.
\ii The dual basis consists of the coordinate rows with a single $1$ in position $j$, because applying the $j$th row to $e_k$ gives $\delta_{jk}$.\end{steps}\end{example*}
\method{Compute a nonstandard dual basis}\begin{steps}\ii In $\mathbb R^2$, use $v_1=(1,1),v_2=(1,-1)$. Solve $(x,y)=av_1+bv_2$ to obtain $a=(x+y)/2$, $b=(x-y)/2$.\ii The coordinate functionals are $\varphi_1(x,y)=(x+y)/2$ and $\varphi_2(x,y)=(x-y)/2$. Evaluating on both basis vectors gives the four required delta values.\ii Equivalently, if $Q$ has columns $v_1,v_2$ in the standard basis, the rows of $Q^{-1}$ are the dual functionals, because $Q^{-1}Q=I$.\end{steps}
\begin{example*}[Notes Chapter 10: a choice-dependent identification]Equal finite dimensions give $V\cong V^{\prime}$, but sending a basis to its dual depends on the chosen basis.
\begin{steps}\ii For $V=\mathbb R$ with basis $(1)$, the dual basis is $\varphi(x)=x$. Sending $1$ to $\varphi$ sends $v$ to the functional $x\mapsto vx$.
\ii With basis $(2)$, the dual basis is $\psi(x)=x/2$. Sending $2$ to $\psi$ sends $v=(v/2)2$ to $(v/2)\psi$, namely $x\mapsto vx/4$.
\ii The two maps differ, for instance at $v=1$. Thus this construction requires a basis choice.\end{steps}\end{example*}
\begin{shortcuts}\ii Dual basis functionals return the original basis coefficients, so solving coordinates computes the dual basis.\ii To prove an identity of functionals, evaluate both sides on a basis; equality there implies equality everywhere.\end{shortcuts}\begin{trapbox}\ii The dual basis does not generally have the same coordinate rows as the original basis columns: $(1,1),(1,-1)$ gives rows $(1/2,1/2),(1/2,-1/2)$.\ii The prime in $V^{\prime}$ does not mean differentiation; a functional such as $p\mapsto p(2)$ returns a scalar, whereas differentiation returns a polynomial.\end{trapbox}
\section{Double duals and separation}\begin{qbox}Do all linear functional values distinguish vectors? For $e_1=(1,0)$, the functional $(x,y)\mapsto x$ takes value $1$, distinguishing $e_1$ from zero.\end{qbox}
\begin{proposition*}[Basis extension and separation tools for Notes Propositions 46, 50--51]Every linear functional on a subspace $U\subseteq V$ extends to one on $V$. If $v\notin U$, there is $\varphi\in V^{\prime}$ with $\varphi(U)=\{0\}$ and $\varphi(v)=1$. For arbitrary infinite-dimensional spaces these statements use the basis-extension form of choice cited in Week 2.\end{proposition*}
\begin{motive}\ii[Want.] Construct a functional detecting a direction outside a subspace.\ii[Know.] A basis of a subspace can be extended to a basis of the full space.\ii[Key step.] Prescribe values on the original basis and arbitrary values on added vectors.\end{motive}
\begin{prob}\begin{proof}\begin{steps}\ii Choose a basis of $U$ and extend it to a basis of $V$. Prescribe the given functional values on the first basis and zero on all added basis vectors.\ii For each finite basis expansion, define the new functional by the same coefficient-weighted sum of prescribed values. Unique expansions make it well-defined; adding and scaling expansions proves linearity. Its restriction equals the original functional.\ii If $v\notin U$, adjoining $v$ to a basis of $U$ is independent: a relation with nonzero coefficient on $v$ would put $v$ in $U$. Extend that list to a basis of $V$.\ii Prescribe $0$ on the basis of $U$, $1$ on $v$, and $0$ on the other basis vectors. The preceding construction gives the separating functional.\end{steps}\end{proof}\end{prob}
\begin{prob}\begin{definition*}[Notes Chapter 10: double dual]The double dual is $V^{\prime\prime}=(V^{\prime})^{\prime}$. Define $j_V(v):V^{\prime}\to F$ by $j_V(v)(\varphi)=\varphi(v)$, and $j_V:V\to V^{\prime\prime}$ by $v\mapsto j_V(v)$.\end{definition*}\end{prob}
\begin{proposition*}[Notes Chapter 10: linearity of evaluation]
For every $v\in V$, $j_V(v)\in V^{\prime\prime}$, and the map $j_V:V\to V^{\prime\prime}$ is linear.
\end{proposition*}
\begin{motive}\ii[Want.] Check that the evaluation rule defines a linear map into the double dual.\ii[Know.] Functionals add pointwise and are linear on vectors.\ii[Key step.] Test linearity first in the functional and then in the vector.\end{motive}
\begin{prob}\begin{proof}\begin{steps}
\ii For fixed $v$, $j_V(v)(a\varphi+\psi)=(a\varphi+\psi)(v)=a\varphi(v)+\psi(v)$. Thus $j_V(v)$ is a linear functional on $V^{\prime}$.
\ii For every $\varphi$, $j_V(av+w)(\varphi)=\varphi(av+w)=a\varphi(v)+\varphi(w)=(aj_V(v)+j_V(w))(\varphi)$. Equality on every functional proves linearity of $j_V$.
\end{steps}\end{proof}\end{prob}
\begin{proposition*}[Notes Proposition 46; LADR 3F.32]The map $j_V:V\to V^{\prime\prime}$ is injective for every vector space. If $V$ is finite-dimensional, it is an isomorphism.\end{proposition*}
\begin{motive}\ii[Want.] Recover a vector from its complete collection of functional values.\ii[Know.] The separation tool supplies a functional nonzero on any nonzero vector.\ii[Key step.] Test whether the image functional $j_V(v)$ can be zero.\end{motive}
\begin{prob}\begin{proof}\begin{steps}\ii For $v\ne0$, apply separation with $U=\{0\}$ to obtain $\varphi(v)=1$. Then $j_V(v)(\varphi)=1$, so $j_V(v)\ne0$. Thus $j_V$ has zero kernel and is injective.\ii If $V$ is finite-dimensional, $\dim V^{\prime\prime}=\dim V^{\prime}=\dim V$. Week 4 Proposition 31 makes this injective map invertible.\ii Consequently if $\varphi(v)=\varphi(w)$ for all $\varphi\in V^{\prime}$, then $j_V(v-w)=0$, and injectivity gives $v=w$.\end{steps}\end{proof}\end{prob}
\begin{shortcuts}\ii To prove $v=0$, it suffices to prove $\varphi(v)=0$ for every linear functional, by separation.\ii To construct a functional with constraints, choose a basis containing the constrained independent vectors and prescribe its values there.\end{shortcuts}\begin{trapbox}\ii One functional usually cannot distinguish every vector: on $\mathbb R^2$, $(x,y)\mapsto x$ cannot distinguish $(0,0)$ from $(0,1)$.\ii $V\cong V^{\prime}$ from a basis choice is different from the specified map $j_V:V\to V^{\prime\prime}$: the former varied in the one-dimensional example, while $j_V(v)(\varphi)=\varphi(v)$ contains no basis choice.\end{trapbox}
\section{Dual maps, transposes and observables}\begin{qbox}How does a functional on the output give a functional on the input? For $T(t)=2t$ and $\varphi(s)=3s$, the input functional is $(\varphi\circ T)(t)=6t$.\end{qbox}
\begin{prob}\begin{definition*}[Notes Definition 28; LADR 3.118]For $T:V\to W$ linear, the dual map is $T^{\prime}:W^{\prime}\to V^{\prime}$, $T^{\prime}(\varphi)=\varphi\circ T$. The composition is a linear functional by Week 3. The domain and codomain are reversed.\end{definition*}\end{prob}
\begin{proposition*}[Notes Proposition 47; LADR 3.120]The map $T^{\prime}$ is linear. The assignment $\mathcal D:L(V,W)\to L(W^{\prime},V^{\prime})$, $\mathcal D(T)=T^{\prime}$, is linear, and for $T:U\to V$, $S:V\to W$, $(ST)^{\prime}=T^{\prime}S^{\prime}$.\end{proposition*}
\begin{motive}\ii[Want.] Calculate dual maps of combinations and compositions.\ii[Know.] Functionals and the original maps are linear; composition is associative.\ii[Key step.] Evaluate both proposed functionals on an arbitrary input vector.\end{motive}
\begin{prob}\begin{proof}\begin{steps}\ii For $\varphi,\psi\in W^{\prime}$, $T^{\prime}(a\varphi+\psi)(v)=(a\varphi+\psi)(Tv)=a\varphi(Tv)+\psi(Tv)$. Thus $T^{\prime}$ is linear.\ii For $T_1,T_2:V\to W$, $((aT_1+T_2)^{\prime}\varphi)(v)=\varphi(aT_1v+T_2v)=a\varphi(T_1v)+\varphi(T_2v)$. This proves $\mathcal D(aT_1+T_2)=a\mathcal D(T_1)+\mathcal D(T_2)$.\ii For $\varphi\in W^{\prime}$, $(ST)^{\prime}\varphi=\varphi\circ S\circ T=T^{\prime}(S^{\prime}\varphi)$, giving the reversed product.\end{steps}\end{proof}\end{prob}
\begin{example*}[Notes Example 45; LADR 3.119]\begin{steps}\ii Let $Tp=xp$ on $\mathcal P(\mathbb R)$ and $\varphi(p)=p(2)$. Then $(T^{\prime}\varphi)(p)=(xp)(2)=2p(2)$, so $T^{\prime}\varphi=2\varphi$.
\ii For differentiation $Dp=p^{\prime}$ and $\psi(p)=p(3)$, $(D^{\prime}\psi)(p)=p^{\prime}(3)$.
\ii If $\eta(p)=\int_0^1p(t)\,dt$, then $(D^{\prime}\eta)(p)=\int_0^1p^{\prime}(t)\,dt=p(1)-p(0)$ by termwise integration.\end{steps}\end{example*}
\begin{prob}\textbf{Dualizing the assignment itself (Notes Chapter 10).}\begin{steps}
\ii For $\mathcal D:T\mapsto T^{\prime}$ above, its dual has domain $L(W^{\prime},V^{\prime})^{\prime}$ and codomain $L(V,W)^{\prime}$ by the definition of a dual map.
\ii That definition gives $\mathcal D^{\prime}(\Lambda)=\Lambda\circ\mathcal D$, hence $(\mathcal D^{\prime}\Lambda)(T)=\Lambda(T^{\prime})$ for each $T$.
\end{steps}\end{prob}
\begin{proposition*}[Notes Proposition 48: double-dual compatibility]For every linear map $T:V\to W$, $T^{\prime\prime}j_V=j_WT$, where $T^{\prime\prime}=(T^{\prime})^{\prime}:V^{\prime\prime}\to W^{\prime\prime}$. This compatibility, together with its basis-free formula, is why $j_V$ is called canonical.\end{proposition*}
\begin{motive}\ii[Want.] Check that applying a map agrees with the double-dual identification.\ii[Know.] $j_V(v)$ evaluates a functional at $v$, and the dual composes a functional with its original map.\ii[Key step.] Evaluate the two double-dual vectors on an arbitrary $\psi\in W^{\prime}$.\end{motive}
\begin{prob}\begin{proof}\begin{steps}\ii For $v\in V$ and $\psi\in W^{\prime}$, $(T^{\prime\prime}j_V(v))(\psi)=j_V(v)(T^{\prime}\psi)=(T^{\prime}\psi)(v)=\psi(Tv)$.\ii The last expression is $j_W(Tv)(\psi)$. Since this holds for every $\psi$, the two elements of $W^{\prime\prime}$ agree, proving the identity.\ii In finite dimension, the $j$ maps are isomorphisms by Proposition 46, so this identity identifies the double-dual map with $T$ through those specified maps.\end{steps}\end{proof}\end{prob}
\begin{proposition*}[Notes Proposition 49; LADR 3.132]For finite-dimensional $V,W$, if $A=M_C^B(T)$, then $M_{B^{\prime}}^{C^{\prime}}(T^{\prime})=A^{\mathsf T}$, where $B^{\prime},C^{\prime}$ are the dual bases. The transpose, defined in Week 3, exchanges rows and columns.\end{proposition*}
\begin{motive}\ii[Want.] Calculate the matrix of a dual map from the original matrix.\ii[Know.] A matrix column contains basis-image coefficients; dual basis evaluations recover those coefficients.\ii[Key step.] Evaluate $T^{\prime}$ of one codomain dual basis vector at each domain basis vector.\end{motive}
\begin{prob}\begin{proof}\begin{steps}\ii Write $B=(v_1,\ldots,v_n)$, $C=(w_1,\ldots,w_m)$, with dual bases $\varphi_k,\psi_j$. Then $Tv_k=\sum_i a_{ik}w_i$.\ii $(T^{\prime}\psi_j)(v_k)=\psi_j(Tv_k)=a_{jk}$. By the functional expansion, $T^{\prime}\psi_j=\sum_k a_{jk}\varphi_k$.\ii Thus entry $(k,j)$ of the dual matrix is $a_{jk}$, exactly entry $(k,j)$ of $A^{\mathsf T}$.\end{steps}\end{proof}\end{prob}
\begin{example*}[Notes Example 46]Let $A=\begin{pmatrix}1&2&3\\4&5&6\end{pmatrix}$ represent $\mathbb R^3\to\mathbb R^2$.
\begin{steps}\ii A functional on $\mathbb R^2$ is a row $(x\ y)$. Precomposing with $A$ gives $(x\ y)A=(x+4y,2x+5y,3x+6y)$.
\ii In dual coordinate columns, the same map is $\begin{pmatrix}1&4\\2&5\\3&6\end{pmatrix}\begin{pmatrix}x\\y\end{pmatrix}$, confirming the transpose formula.\end{steps}\end{example*}
\begin{proposition*}[Notes Chapter 10: distributions and observables]For a column-stochastic matrix $P$, let a probability column $p$ describe the present state and let the real column $f$ list rewards $f_i$ at each state. Their mean reward is $f^{\mathsf T}p=\sum_i f_ip_i$. After one step it is $f^{\mathsf T}Pp=(P^{\mathsf T}f)^{\mathsf T}p$. The value $(P^{\mathsf T}f)_j=\sum_iP_{ij}f_i$ is the conditional mean reward when the present state is $j$.\end{proposition*}
\begin{motive}\ii[Want.] Calculate the same expected reward by changing the distribution or the reward function.\ii[Know.] Column $j$ of $P$ contains the next-state probabilities from state $j$ (Week 3).\ii[Key step.] Regroup the finite weighted sum by present state.\end{motive}
\begin{prob}\begin{proof}\begin{steps}\ii After one step the distribution is $Pp$, so the mean is $\sum_i f_i\sum_jP_{ij}p_j=\sum_j(\sum_iP_{ij}f_i)p_j$.\ii The inner sum is the mean of the rewards weighted by the conditional probabilities in column $j$. These inner sums are exactly the entries of $P^{\mathsf T}f$, proving the identity.\ii For the row convention, put $A=P^{\mathsf T}$ and initial distribution row $\lambda^{\mathsf T}$. For integers $k\ge0$, with $A^0=I$ and $A^k$ the repeated product defined in Week 3, $\lambda^{\mathsf T}A^kf$ is the mean after $k$ steps, by repeated multiplication and transposition.\end{steps}\end{proof}\end{prob}
\method{Compute an expected reward in two ways}\begin{steps}\ii Use Week 3's $P=\left(\begin{smallmatrix}1/2&1/4&0\\1/2&1/2&1/2\\0&1/4&1/2\end{smallmatrix}\right)$, initial $p=e_1$, and rewards $f=(0,1,4)^{\mathsf T}$. Then $Pp=(1/2,1/2,0)^{\mathsf T}$ and $f^{\mathsf T}Pp=1/2$.\ii Alternatively $P^{\mathsf T}f=(1/2,3/2,5/2)^{\mathsf T}$, whose first entry gives the same mean from state 1.\end{steps}
\begin{shortcuts}\ii Compute a dual map by precomposition; its matrix in dual bases is the transpose, so no new matrix solving is needed.\ii For expected rewards, choose whichever is shorter: update $p$ by $P$ or update $f$ by $P^{\mathsf T}$.\end{shortcuts}\begin{trapbox}\ii Dual maps reverse direction: a $2\times3$ matrix gives a dual matrix of size $3\times2$, not $2\times3$.\ii Transposition does not reverse entries within a row; it exchanges indices. Example 46 explicitly gives the correct array.\ii Unreversed dual products may be undefined: for $T:\mathbb R\to\mathbb R^2$ and $S:\mathbb R^2\to\mathbb R^3$, $T^{\prime}S^{\prime}$ is defined, but $S^{\prime}T^{\prime}$ has mismatched intermediate spaces.\end{trapbox}
\section{Annihilators and rank of the dual map}\begin{qbox}Which functionals vanish on a specified subspace? For the plane spanned by $(1,1,0),(0,1,1)$, the row $(-1,1,-1)$ vanishes on both generators, hence on the whole plane.\end{qbox}
\begin{prob}\begin{definition*}[Notes Definition 29; LADR 3.121]The annihilator of a subspace $U\subseteq V$ is $U^0=\{\varphi\in V^{\prime}:\varphi(u)=0\text{ for every }u\in U\}$. It is a set of functionals in $V^{\prime}$, not of vectors in $V$.\end{definition*}\end{prob}
\begin{proposition*}[LADR 3.124--3.125, 3.127]$U^0$ is a subspace of $V^{\prime}$. In finite dimension, $\dim U^0=\dim V-\dim U$. In all dimensions, $U^0=\{0\}$ iff $U=V$, and $U^0=V^{\prime}$ iff $U=\{0\}$.\end{proposition*}
\begin{motive}\ii[Want.] Count functionals that vanish on $U$ and detect the extreme cases.\ii[Know.] Dual basis expansions, basis extension, and the separation tool are available.\ii[Key step.] Extend a basis of $U$; the dual coefficients on that first basis must vanish.\end{motive}
\begin{prob}\begin{proof}\begin{steps}\ii The zero functional annihilates $U$, and if $\varphi,\psi$ do, then $(a\varphi+\psi)(u)=0$ for every $u\in U$. Thus $U^0$ is a subspace.\ii In finite dimension, extend a basis $u_1,\ldots,u_k$ of $U$ to $u_1,\ldots,u_n$ in $V$ and take its dual basis $\varphi_1,\ldots,\varphi_n$. A functional annihilates $U$ iff its first $k$ coefficients in this dual basis are zero. Thus $\varphi_{k+1},\ldots,\varphi_n$ is a basis of $U^0$, of length $n-k$.\ii If $U=V$, only the zero functional vanishes everywhere. If $U\ne V$, separation supplies a nonzero functional annihilating $U$, so $U^0\ne\{0\}$.\ii If $U=\{0\}$, every functional vanishes there. If $U$ contains a nonzero $u$, separation supplies a functional with value $1$ at $u$, so $U^0\ne V^{\prime}$.\end{steps}\end{proof}\end{prob}
\begin{example*}[Notes Example 47]Let $U=\operatorname{span}((1,1,0),(0,1,1))\subseteq\mathbb Q^3$.
\begin{steps}\ii A row $(a,b,c)$ annihilates both generators iff $a+b=0$ and $b+c=0$.
\ii Thus $(a,b,c)=b(-1,1,-1)$, so $U^0=\operatorname{span}((-1,1,-1))$.
\ii The two generators are independent by their first and last coordinates, so $\dim U=2$ and $\dim U^0=1$, verifying $2+1=3$.\end{steps}\end{example*}
\begin{proposition*}[Notes Proposition 50; LADR 3.128(a), 3.130(b)]For any linear $T:V\to W$, $\ker T^{\prime}=(\im T)^0$ and $\im T^{\prime}=(\ker T)^0$. The general-dimensional statement uses the functional-extension tool above.\end{proposition*}
\begin{motive}\ii[Want.] Find the kernel and range of the dual from the original map.\ii[Know.] $T^{\prime}\psi=\psi\circ T$, so vanishing of this composite is a condition on $\im T$.\ii[Key step.] For the range converse, first define a functional on $\im T$, then extend it to $W$.\end{motive}
\begin{prob}\begin{proof}\begin{steps}\ii $\psi\in\ker T^{\prime}$ iff $\psi(Tv)=0$ for all $v$, exactly $\psi\in(\im T)^0$.\ii If $\varphi=T^{\prime}\psi$ and $v\in\ker T$, then $\varphi(v)=\psi(Tv)=\psi(0)=0$. Thus $\im T^{\prime}\subseteq(\ker T)^0$.\ii For $\varphi\in(\ker T)^0$, define $\eta:\im T\to F$ by $\eta(Tv)=\varphi(v)$. If $Tv=Tw$, then $v-w\in\ker T$, so $\varphi(v)=\varphi(w)$; the rule is well-defined.\ii It is linear: $\eta(aTv+Tw)=\eta(T(av+w))=\varphi(av+w)=a\varphi(v)+\varphi(w)$. Extend $\eta$ to $\widetilde\eta\in W^{\prime}$ by the extension tool. Then $T^{\prime}\widetilde\eta(v)=\widetilde\eta(Tv)=\varphi(v)$, proving the reverse inclusion.\end{steps}\end{proof}\end{prob}
\begin{proposition*}[Notes Proposition 51; LADR 3.129, 3.131]For a linear map $T:V\to W$, $T$ is injective iff $T^{\prime}$ is surjective; $T$ is surjective iff $T^{\prime}$ is injective.\end{proposition*}
\begin{motive}\ii[Want.] Determine injectivity or surjectivity from the dual.\ii[Know.] Proposition 50 expresses the dual kernel and range as annihilators; the preceding extreme-case criteria apply in all dimensions.\ii[Key step.] Substitute zero kernel or full range into those identities.\end{motive}
\begin{prob}\begin{proof}\begin{steps}\ii $T$ injective $\iff\ker T=\{0\}\iff(\ker T)^0=V^{\prime}\iff\im T^{\prime}=V^{\prime}\iff T^{\prime}$ surjective.\ii $T$ surjective $\iff\im T=W\iff(\im T)^0=\{0\}\iff\ker T^{\prime}=\{0\}\iff T^{\prime}$ injective.\end{steps}\end{proof}\end{prob}
\begin{proposition*}[LADR 3.128(b), 3.130(a), 3.133]For a linear map $T:V\to W$ between finite-dimensional spaces, $\operatorname{rank}T^{\prime}=\operatorname{rank}T$ and $\dim\ker T^{\prime}=\dim\ker T+\dim W-\dim V$. Consequently row rank equals column rank for every finite matrix.\end{proposition*}
\begin{motive}\ii[Want.] Translate duality identities into dimension formulas.\ii[Know.] The annihilator dimension formula and rank--nullity.\ii[Key step.] Count the annihilators in Proposition 50.\end{motive}
\begin{prob}\begin{proof}\begin{steps}\ii $\dim\im T^{\prime}=\dim(\ker T)^0=\dim V-\dim\ker T=\operatorname{rank}T$.\ii $\dim\ker T^{\prime}=\dim(\im T)^0=\dim W-\operatorname{rank}T=\dim W-\dim V+\dim\ker T$.\ii For a matrix $A$, its column rank is the rank of its map by Week 3. The dual matrix is $A^{\mathsf T}$, whose column rank is the row rank of $A$. Equality of map ranks gives the equality of matrix ranks, confirming Week 3's factorization proof.\end{steps}\end{proof}\end{prob}
\begin{shortcuts}\ii To compute an annihilator, write a general row and set its value on each subspace basis vector equal to zero.\ii Finite-dimensional annihilator dimensions supply a quick check: subspace dimension plus annihilator dimension equals ambient dimension.\end{shortcuts}\begin{trapbox}\ii Annihilator containment reverses subspace containment: more vectors to vanish on means fewer functionals. For $0\subset\mathbb Re_1\subset\mathbb R^2$, the annihilators are all rows, rows $(0,b)$, and zero.\ii Equal rank of $T$ and $T^{\prime}$ does not give equal nullity when domain and codomain dimensions differ: the map $\mathbb R^2\to\mathbb R$, $(x,y)\mapsto x$, has nullity $1$, but its dual has nullity $0$.\end{trapbox}
\section{Problem set}\begin{qbox}Use quotient representatives, dual basis values, and annihilator equations to reduce each problem to a map already defined.\end{qbox}
\begin{prob}\textbf{HW5 P1.} Let $T:\mathcal P_3(\mathbb R)\to\mathbb R^2$, $Tp=(p(0),p(1))$. (a) Determine $\ker T$ and prove surjectivity. (b) Describe the induced isomorphism $\mathcal P_3(\mathbb R)/\ker T\to\mathbb R^2$. (c) Find a quotient basis.\\\emph{Hint:} Factor polynomials vanishing at 0 and 1; use linear polynomials for arbitrary outputs.\end{prob}
\begin{prob}\textbf{HW5 P2.} If $T:V\to W$ is linear and $U$ is a subspace of $V$ with $U\subseteq\ker T$, prove there is a unique linear $\widetilde T:V/U\to W$ with $\widetilde T(v+U)=Tv$. Determine its kernel.\\\emph{Hint:} Use Proposition 42 with target subspace zero, then test the zero output.\end{prob}
\begin{prob}\textbf{HW5 P3 (LADR 3F.1).} Prove every linear functional $\varphi:V\to F$ is either zero or surjective.\\\emph{Hint:} Scale a vector where the value is nonzero.\end{prob}
\begin{prob}\textbf{HW5 P4.} For a basis $v_1,\ldots,v_n$ and its dual basis $\varphi_1,\ldots,\varphi_n$, prove (a) $v=\sum_j\varphi_j(v)v_j$ for each $v\in V$ and (b) $\psi=\sum_j\psi(v_j)\varphi_j$ for each $\psi\in V^{\prime}$.\\\emph{Hint:} Apply a dual functional to a basis expansion, then evaluate the second identity on each $v_k$.\end{prob}
\begin{prob}\textbf{HW5 P5.} Let $T:V\to W$ be linear between finite-dimensional spaces and let $A$ represent it in bases $B,C$. Prove $T^{\prime}:W^{\prime}\to V^{\prime}$ has matrix $A^{\mathsf T}$ in the corresponding dual bases.\\\emph{Hint:} Compute $(T^{\prime}\psi_i)(v_j)$, as in Proposition 49.\end{prob}
\begin{prob}\textbf{HW5 P6.} Let $U=\operatorname{span}((1,1,0),(0,1,1))\subseteq\mathbb R^3$. (a) Find $U^0$, identifying functionals with rows. (b) Verify $\dim U+\dim U^0=3$.\\\emph{Hint:} A general row must vanish on both generators.\end{prob}
\begin{prob}\textbf{HW5 P7 (optional).} For $j_V(v)(\varphi)=\varphi(v)$, prove $j_V:V\to V^{\prime\prime}$ is injective, and is an isomorphism when $V$ is finite-dimensional.\\\emph{Hint:} Use separation to detect a nonzero vector and the equality of dual dimensions.\end{prob}
\begin{prob}\textbf{Exam practice (practice-midterm2.pdf, Question 2).} Let $V,W,Z$ be finite-dimensional real vector spaces, and let $T:V\to W$, $S:W\to Z$ be linear. (a) Define the dual map $T^*:W^*\to V^*$ and prove it is linear. (b) Prove $(S\circ T)^*=T^*\circ S^*$. (c) Prove that if $T$ is surjective, then $T^*$ is injective. (d) For $V=\mathbb R^2$, $W=\mathbb R$, $T(x,y)=2x-y$, and $\ell(t)=3t$, compute $T^*\ell$. In this exam, $*$ denotes the same dual-space and dual-map notation as our prime.\\\emph{Hint:} Use precomposition for (a,b), and evaluate a functional on all outputs of $T$ for (c).\end{prob}\end{document}
